\section{The local theorems and their proof routes}
\label{sec:v35-literature-routes}
There are two local questions in this article. The stationary physical
question is a lattice singleton in three coordinates, evaluated at a
deterministic physical time. The actual-return question is a
four-coordinate mixed density with three counting coordinates and a
pointwise roof coordinate. They have different regularizing mechanisms.
The first uses the actual stationary flight age; the second requires
the extracted raw return correction. We keep both records, rather than
identify one with the other.

\subsection{Comparison at the level of theorem statements}
The Lorentz-process local theorem of Sz\'asz--Varj\'u \cite{SV} concerns
the periodic cell process and its recurrence, using a Young-tower local
limit framework. Its existence already prevents interpreting a Gaussian
cell local limit at a fixed table as a new phenomenon here. The endpoint
version in \cite[Theorem 2.2 and Remark 2.3]{DPZ} gives a mixing local
limit for the two-dimensional cell index with piecewise H\"older
endpoint observables, including in randomly deformed configurations.
Thus endpoint factors and geometric perturbations are not by themselves
new local-limit principles. That theorem is not, as stated, the joint
three-frequency displacement--roof estimate used in this article.
The new work required for that use is the accumulated roof action in
all three norms and its uniform compact-frequency arithmetic.

The suspension theory of Dolgopyat--N\'andori
\cite[Definition 2.2, Theorem 3.7 and Proposition 6.3]{DN} starts from
a mixing local limit and the appropriate closed arithmetic group, and
includes finite-horizon Sinai billiards among its applications. At a
fixed parameter, the physical conclusion here belongs to that
framework once the joint observable, its endpoint coboundaries and
its minimality hypotheses have been verified. The scalar billiard
application there is not a literal statement of our compact-family
vector theorem, but neither is the suspension step claimed as a new
mechanism. Our proof supplies the needed collision input directly and
tracks its uniformity through the actual age integration.

The collision count is also not a new arithmetic group merely because
it is recorded at a continuous time. If $x_s$ is the last outgoing
collision state at time $s$, the starting age is $a$, the ending state
is $y=T^mx$, and the ending age is $v=a+t-S_m$, then exactly
\begin{equation}\label{eq:v35-count-as-suspension-reward}
 \int_0^t\frac{ds}{\tau(x_s)}
                     =m-\frac{a}{\tau(x)}+\frac{v}{\tau(y)}.
\end{equation}
The incomplete first and last flights prove this identity, with the
same formula when $m=0$. Thus the count is an additive suspension
reward up to a bounded endpoint coboundary. The arithmetic to check
is the joint lattice/roof group with these endpoints. Our finite-cover
argument performs that check without a rotation of the obstacle.

There is no claim that a singleton in a discrete lattice coordinate
is a stronger topology than a mixing lattice local limit: a singleton
indicator is continuous on the lattice. Nor is a fixed roof interval
silently replaced by a pointwise roof density. The positive age overlap
turns the physical event into an admissible interval-type endpoint
test. That regularization is absent for the raw actual-return density.
The source total-variation conclusion has a different purpose: after
a positive microscopic denominator is proved, it controls every
bounded test on the unchanged conditioned trajectory. It does not
supply an arbitrary selector's Gaussian amplitude.

The contribution developed in this part is consequently a specific
combination of verifiable statements: an action-weighted norm argument
for the unsmoothed roof, a rotation-free finite-cover exclusion of
joint phases, a compact-family collision-to-physical local theorem,
and a quantitative source comparison after normalization. The ellipse
family verifies these ingredients beyond the circular parameter line.
No claim of priority for general local inversion, abstract suspension
local limits, or perturbations of dispersing tables is needed for these
statements.

\subsection{A short route through the article}
The leading physical theorem can be checked in the following order.
Section~\ref{sec:action-norm-details} specifies the exact norms,
partition conditions, Jacobian sums, action estimates and physical
peripheral representation. Section~\ref{sec:rotation-free-arithmetic}
excludes the joint phase by contact holonomy and finite-cover mixing.
Section~\ref{sec:compact-family-local-principle} derives the compact-family
local law, verifies the ellipse geometry, and proves the normalized
source comparison. Its only imported dynamical inputs are the
unweighted finite-horizon collision framework of \cite{DZ11,DZ}, the
piecewise multiplier criterion of \cite{DPZ}, and the product/mixing
construction of \cite{Young}.

Sections~\ref{sec:compact-collision-spectrum}--\ref{sec:stationary-microscopic-LLT}
retain the detailed circular collision, endpoint-measure and age-overlap
proof. The exact physical identity can also be checked directly in
Section~\ref{sec:exact-stationary-endpoint}; the positive likelihood
there is used in Corollary~\ref{cor:v35-family-posterior}.
The signed-complement conclusion additionally uses the exact
Gaussian-plus-complement identity in
Section~\ref{sec:stationary-gaussian-complement}; it is a consequence
of the independently evaluated positive probability, not a uniform
spectral estimate at a growing band.

Part II retains the full actual-return analysis. Its independent
unconditional results include the moving-section geometry, exponential
return bounds, Gaussian and functional laws, complete joint
nondegeneracy and marked moments. Its raw inversion route consists of
the critical-edge calculation, finite-count extraction, coherent cutoff
transport, and the remaining common pointwise correction and return
complement. In particular Theorem~\ref{thm:LLT} keeps its original
four-coordinate object and its exact analytic hypotheses. The new
three-coordinate theorem is not substituted into that statement.
The A--X synopsis remains compiled in
Appendix~\ref{app:retained-return-statements}. The original sources
and their proof routes remain available even where the direct physical
argument bypasses them.
